\subsection{Solving Equations: Basic Moves} 
The properties of equality give us two basic moves that we can use to solve equations.
\begin{itemize}
\item We can \makeblank[1in] the same number to each side of an equation.\lbr We can also \makeblank[1.5in] (this is just adding the opposite).
\item We can \makeblank[1.5in] by the same number on each side of an equation.\lbr We can also \makeblank[1.5in] (this is just multiplying by the inverse).
\end{itemize}
\begin{Note}
We can \makeblank[1in] any number, so it is safe to \makeblank[1in] variable expressions. However, we must be careful not to \makeblank[2in] by zero, so we don't \makeblank[2in] by variable expressions at this stage.
\end{Note}
We choose basic moves in order to isolate the variable on one side.
\begin{Example}
Solve $2x-3=7$
\why{4}
\end{Example}
\begin{Example}
Solve $5x-7=2x+5$
\why{5}
\end{Example}
\begin{Example}
Solve $\frac{2}{3}x-\frac{3}{5}=\frac{1}{6}x+\frac{4}{5}$
\why{7}
\end{Example}